% !TeX root = ../../Book2.tex
\section{完备性}
\label{b2:sec:1003}

两两正交的模态未必能恢复整个空间。本节用移位逆的紧性排除遗漏的正交方向，再区分普通平方范数、一次能量权和两次算子权；求解与求导分别需要不同的收敛条件。

\subsection{紧算子方法}
\label{b2:subsec:100301}

算子 \(A\) 对高频求导，通常无界；其移位逆 \(K\) 则是具有紧性的积分算子。若实或复 Banach 空间 \(X,Y\) 间的有界线性算子 \(T:X\to Y\) 把每个有界列送到有收敛子列的像列，则称它为\term[jin-suanzi]{紧算子}[compact operator]。这里使用第一卷定理 8.5.5 的紧自伴谱定理：紧自伴算子的非零特征向量张成核空间的正交补，每个非零特征值重数有限，唯一可能的聚点是零。应用前须核对紧性及核，而不能只从自伴性猜测完备性。

\Needspace{6\baselineskip}
\begin{theorem}[Sturm--Liouville 展开]\label{b2:thm:10-sl-spectral}
设 \(I=(a,b)\) 为有限区间，\(p,w,q\in C([a,b];\mathbb R)\)，\(p,w>0\)；\(Au=w^{-1}(-(pu')'+qu)\) 取两端各一个实分离 Dirichlet 或 Robin 条件，定义域如\cref{b2:def:10-sl-operator}，\(H=L^2(I,w\,\mathrm dx)\)。令 \(V\subset H^1(I)\) 由 Dirichlet 端点值为零限定，\(R\) 为 Robin 端点集合，\(h_s\in\mathbb R\)，\(\mathfrak a(u,v)=\int_I(pu'\overline{v'}+qu\overline v)\,\mathrm dx+\sum_{s\in R}h_su(s)\overline{v(s)}\)。取 \(\sigma>0\)，使移位形式强制且 \(\mathfrak a(u,u)+\sigma\|u\|_H^2\geq\|u\|_H^2\)。存在一组实的标准正交特征函数 \((e_n)_{n\geq1}\)，构成 \(H\) 的基。相应特征值可按严格递增次序排列为
\[
 -\sigma<\lambda_1<\lambda_2<\cdots,
 \quad\lambda_n\longrightarrow+\infty.
\]
特别地，负特征值只有有限多个。每个 \(f\in H\) 都有
\[
 f=\sum_{n=1}^{\infty}f_ne_n,
 \quad f_n=\langle f,e_n\rangle_H,
 \quad\|f\|_H^2=\sum_{n=1}^{\infty}|f_n|^2,
\]
其中级数在 \(H\) 中收敛。
\end{theorem}
\begin{proof}
将 \([a,b]^2\) 分成有限个小矩形，在每个矩形内以 \(G_\sigma\) 的一个取值近似它。连续性给出阶梯核 \(G_j\to G_\sigma\) 一致收敛。对应算子 \(K_j\) 的像落在有限个区间示性函数的张成空间内，因而有限秩。记 \(m_w=\int_a^bw\)，则
\[
 \|(K-K_j)f\|_H
 \leq m_w\|G_\sigma-G_j\|_\infty\cdot\|f\|_H.
\]
故 \(K_j\to K\) 于算子范数。具体地，给定单位球像的覆盖精度 \(\varepsilon\)，先取 \(j\) 使 \(\|K-K_j\|<\varepsilon/3\)。有限维空间中的有界集 \(K_j(B_H)\) 有有限个半径 \(\varepsilon/3\) 的覆盖球，同样的球心以半径 \(2\varepsilon/3\) 覆盖 \(K(B_H)\)。所以后者全有界；\(H\) 完备使其闭包紧，得到 \(K\) 紧。第一卷紧自伴谱定理与\cref{b2:prop:10-green-resolvent}给出正特征值 \(\mu_n\) 及其标准正交特征向量；因 \(\ker K=0\)，这些向量在全部 \(H\) 中完备。这个列必须无穷，因为 \(H\) 无穷维，而有限秩算子不能在无穷维空间上单射。因此 \(\mu_n\to0\)。

若 \(Ke_n=\mu_ne_n\)，则 \(e_n\in\operatorname{ran}K=D(A)\)，且 \(Ae_n=(\mu_n^{-1}-\sigma)e_n\)。反过来，任一 \(A\) 的特征向量都是 \(K\) 的特征向量。上一节的单重性使 \(\mu_n\) 无重复，并可取实特征函数；实系数常微分方程的左端初值先取为实数即可。按 \(\mu_n\) 递减排序，得到所述严格递增的 \(\lambda_n\)。完备性给出展开及 Parseval 等式。
\end{proof}

这个定理从 \(n=1\) 起按特征值递增统一编号。它与\subsecref{b2:subsec:100103}为突出 Neumann 零模态而使用的 \(n=0,1,2,\ldots\) 是同一组函数的两种编号。例如对 \((0,\ell)\)、\(\ell>0\) 上的 \(-u''\) 取两端 Neumann 条件，本节可选
\[
 \lambda_1=0,\quad e_1=\ell^{-1/2},\qquad
 \lambda_{j+1}=\left(\frac{j\pi}{\ell}\right)^2,\quad
 e_{j+1}(x)=\sqrt{\frac2\ell}\cos\frac{j\pi x}{\ell}\quad(j\geq1).
\]
也就是说，本节的 \(e_1\) 对应前面的 \(e_0\)，本节的 \(e_{j+1}\) 对应前面的 \(e_j\)。后面的演化级数沿用本节的统一编号，零模态应按 \(\lambda_n=0\) 识别。

\ProseParagraphBreak

这里的两个谱极限方向相反：\(K\) 的小特征值对应 \(A\) 的大特征值。零始终是 \(K\) 的谱极限，却不是它的特征值；\(A\) 是否有零特征值，则取决于原方程及边界条件。移位只是为证明选择一个可逆点，并没有把物理模型中的零频率删去。

\ProseParagraphBreak

有限区间的假设不能仅从符号上删掉。考虑实线上 \(-u''+u=f\)，其平方可积解算子为与 \(k(x)=e^{-|x|}/2\) 卷积。分布求导给出 \(-k''+k=\delta_0\)，而 \(\widehat k(\xi)=(1+4\pi^2\xi^2)^{-1}\)，所以这是有界单射的解算子。然而它不紧：固定非零 \(f\in C_c^\infty\)，记 \(g=k*f\ne0\)。平移后的 \(f(\,\cdot-jR)\) 范数不变，相应的解为 \(g(\,\cdot-jR)\)。平方可积函数的远平移内积趋零，这可先把 \(g\) 在大紧区间外的平方尾截掉，再使用两个紧支集的分离证明。选足够稀疏的平移序列，相应解之间的距离便有同一个正下界，不可能有收敛子列。

\ProseParagraphBreak

这个反例中，局部积分仍然提高正则性，但函数可以向无穷远整体移动，阻止了全空间紧性。本章连续核在紧正方形上的有限秩逼近同时利用了局部平滑和空间范围有界；后一项在实线模型中消失。因而全空间的 Fourier 积分不能无条件改写成一列平方可积特征函数的展开。

\subsection{特征函数展开}
\label{b2:subsec:100302}

为了使用展开求解方程，需要知道哪些系数列还可以求导。记 \(r_n=\lambda_n+\sigma\geq1\)，并令 \(P_Nf=\sum_{n=1}^Nf_ne_n\)。以下两个刻画区分算子定义域与能量空间。

\Needspace{6\baselineskip}
\begin{proposition}[定义域和能量的系数刻画]\label{b2:prop:10-spectral-domains}
设 \(I=(a,b)\) 为有限区间，\(p,w,q\in C([a,b];\mathbb R)\)，\(p,w>0\)；\(Au=w^{-1}(-(pu')'+qu)\) 取两端各一个实分离 Dirichlet 或 Robin 条件，定义域如\cref{b2:def:10-sl-operator}，\(H=L^2(I,w\,\mathrm dx)\)。令 \(V\subset H^1(I)\) 由 Dirichlet 端点值为零限定，\(R\) 为 Robin 端点集合，\(h_s\in\mathbb R\)，\(\mathfrak a(u,v)=\int_I(pu'\overline{v'}+qu\overline v)\,\mathrm dx+\sum_{s\in R}h_su(s)\overline{v(s)}\)。取强制的正移位 \(\mathfrak a_\sigma=\mathfrak a+\sigma\langle\cdot,\cdot\rangle_H\)，令 \((e_n,\lambda_n)\) 为 \(A\) 的标准正交谱基，\(r_n=\lambda_n+\sigma>0\)、\(f_n=\langle f,e_n\rangle_H\)。对 \(f\in H\)，有
\[
 f\in D(A)\quad\Longleftrightarrow\quad
 \sum_n r_n^2|f_n|^2<\infty,
\]
此时 \(Af=\sum_n\lambda_nf_ne_n\) 在 \(H\) 中成立。又有
\begin{equation}\label{b2:eq:10-form-parseval}
 f\in V\quad\Longleftrightarrow\quad\sum_n r_n|f_n|^2<\infty,
 \quad\mathfrak a_\sigma(f,f)=\sum_n r_n|f_n|^2.
\end{equation}
对应截断分别在定义域范数 \(\|f\|_H+\|Af\|_H\) 与 \(H^1(I)\) 范数中收敛。
\end{proposition}
\begin{proof}
若 \(f\in D(A)\)，令 \(g=(A+\sigma I)f\)，则 \(f=Kg\)，系数关系为 \(g_n=r_nf_n\)。反过来，若 \((r_nf_n)\in\ell^2\)，它给出某个 \(g\in H\)，而 \(Kg\) 与 \(f\) 的全部系数相同，故 \(f=Kg\in D(A)\)。这也证明定义域范数中的截断收敛。

\(V\) 是 \(H^1\) 的闭子空间，因为端点取值连续；由\cref{b2:lem:10-form-coercive}，\(\mathfrak a_\sigma\) 是它的一个等价 Hilbert 内积。函数 \(e_n/\sqrt{r_n}\) 在这个内积下标准正交。若 \(v\in V\) 与它们全部正交，分部积分给出
\[
 0=\mathfrak a_\sigma(v,e_n)=r_n\langle v,e_n\rangle_H,
\]
而 \(H\) 中的完备性推出 \(v=0\)。所以它们也是 \(V\) 的标准正交基，\(V\) 中的 Fourier 系数为 \(\sqrt{r_n}v_n\)，由此得式\eqref{b2:eq:10-form-parseval}的正向及等式。反向中若加权和有限，\(P_Nf\) 在 \(V\) 中为 Cauchy 列，其 \(V\) 极限在 \(H\) 中只能为 \(f\)，故 \(f\in V\)。
\end{proof}

取有限区间 \(I=(a,b)\)、\(p,w,q\in C([a,b];\mathbb R)\)、\(p,w>0\)，令 \(A\) 为\cref{b2:def:10-sl-operator}中具有实分离 Dirichlet 或 Robin 边界条件的算子，\(H=L^2(I,w\,\mathrm dx)\)。取 \(f\in H\)，以 \((e_n)\) 为其标准正交特征基，\(Ae_n=\lambda_ne_n\)、\(f_n=\langle f,e_n\rangle_H\)、\(r_n=\lambda_n+\sigma\)，其中 \(\sigma>0\) 使 \(A+\sigma I\) 正定。系数刻画也给出非齐次边值问题的完整可解条件。若实数 \(z\) 不等于任何 \(\lambda_n\)，则
\[
 u=\sum_n\dfrac{f_n}{\lambda_n-z}e_n
\]
属于 \(D(A)\)，并且 \((A-zI)u=f\)。因为 \(\lambda_n\to\infty\)，比值 \(r_n/|\lambda_n-z|\) 一致有界，定义域条件正由这个界保证。若 \(z=\lambda_k\)，可解当且仅当 \(f_k=0\)；满足这个条件时去掉第 \(k\) 项所得级数给出一个解，全部解还可加上任意倍 \(e_k\)。这里的自由常数来自齐次核。

\ProseParagraphBreak

例如 Neumann 问题 \(-u''=f\) 的可解条件为 \(\int_0^\ell f=0\)，因为零特征函数为常数。将方程直接积分也得到这一条件；谱方法进一步给出它的充分性和除常数外的唯一性。取零均值解相当于在核的正交补中选取唯一代表。

\subsection{Fourier 级数作为特殊情形}
\label{b2:subsec:100303}

对 \(p=w=1,q=0\)，\cref{b2:thm:10-sl-spectral}与显式计算给出正弦基、余弦基及混合边界基的完备性。正弦展开也可以把区间函数作奇延拓后由周期 Fourier 理论得到，余弦展开则对应偶延拓。两条路线各有用途：周期化便于使用卷积核和求和公式，Green 算子路线则能在 \(p,w,q\) 改变后继续成立。

\ProseParagraphBreak

能量刻画把熟悉的导数公式变成精确命题。对 \(f\in H_0^1(0,\ell)\)，其正弦系数满足
\[
 \int_0^\ell|f'|^2\,\mathrm{d}x
 =\sum_{n\geq1}\left(\dfrac{n\pi}{\ell}\right)^2|f_n|^2.
\]
因而
\[
 \|f\|_2\leq\dfrac\ell\pi\|f'\|_2,
\]
等号恰在第一正弦模态及零函数处成立。Neumann 情形对零均值的 \(H^1\) 函数有同一最优常数，但不加零均值条件时常数函数使不等式失败。这些结论来自同一能量恒等式；边界值与均值限制决定了最低允许频率。

\ProseParagraphBreak

一般系数问题仍有 \(H\) 与 \(V\) 中的截断收敛，却不能仅由本章谱定理断言任意连续函数的原始特征函数级数处处收敛。正交完备、点态收敛与一致收敛是不同层次。下一节以范数极限解释演化方程，正是为避免把第一种结论误当成后两种。

\ProseParagraphBreak

可以把这些区别放到一个可算的初始形状中。取 \(I=(0,\pi)\) 的 Dirichlet 条件，令 \(f(x)=\min\{x,\pi-x\}\)。它在中点有折角，弱导数在左右两半分别为 \(1\)、\(-1\)。对非单位正弦系数作一次分部积分，得到
\[
 b_n=\dfrac2\pi\int_0^\pi f(x)\sin(nx)\,\mathrm{d}x
 =\dfrac{4\sin(n\pi/2)}{\pi n^2}.
\]
级数绝对一致收敛，且由平方范数完备性恢复这个连续函数。单位系数为 \(f_n=\sqrt{\pi/2}\,b_n\)，所以
\[
 \sum_n n^2|f_n|^2=\dfrac8\pi\sum_{n\text{ 为奇数}}\dfrac1{n^2}
 =\int_0^\pi|f'|^2\,\mathrm{d}x=\pi.
\]
一次能量权可和，但 \(\sum_n n^4|f_n|^2\) 发散；这个形状属于 \(V\) 而不属于 \(D(A)\)。展开收敛与能否再求二阶导数，因而在同一个对象上有了不同答案。

